% ============================================
% SECTION 7: BEST PRACTICES
% ============================================
\section{Best Practices for FDE Programs}

\subsection{Three Conditions for Success (McCardel Framework)}

Barry McCardel, drawing on five years as a Palantir FDE and subsequent
experience building Hex, identifies three non-negotiable conditions for
FDE program viability:

\begin{enumerate}
    \item \textbf{High-value markets}: FDE economics only work for
        Fortune~500 and government customers. The contract floor is
        seven figures. Below that threshold, the model destroys value.

    \item \textbf{Epistemological humility about product direction}:
        The company must genuinely not know what the product should look
        like for a given market segment, and must use FDEs as the
        discovery mechanism. If the product roadmap is already clear, FDEs
        add cost without insight.

    \item \textbf{Extensible platform architecture}: The platform must be
        designed so that FDE-built solutions can be generalized. If the
        architecture does not support abstraction, every FDE engagement
        produces a dead-end custom build.
\end{enumerate}

\subsection{Platform-First Discipline}

The single most important organizational principle: the company must be a
product company that uses FDEs, not a consulting company that builds a
product. The acid test from McCardel: \textit{``If FDEs prefer one-offs
over your abstraction, it's the wrong abstraction.''}

\subsection{Operational Best Practices}

\begin{itemize}
    \item \textbf{Feedback loop field-to-product}: Every deployment must
        leave reusable assets for the next engagement. If an FDE builds
        something and it dies with the engagement, the program is failing.
    \item \textbf{Clear client boundaries}: FDEs must know when to say
        no. Unlimited scope creep transforms FDEs into outsourced staff.
    \item \textbf{Ethics and compliance from day zero}: FDEs operating
        inside customer environments handle sensitive data and systems.
        Compliance cannot be an afterthought.
    \item \textbf{Rotation cadence}: Prevent burnout by rotating FDEs
        between field deployments and HQ/platform work. Palantir's
        Project Frontline was an early attempt at institutionalizing this.
\end{itemize}

% ============================================
% SECTION 8: RECRUITMENT
% ============================================
\section{Recruiting FDEs: Process \& Pitfalls}

\subsection{Palantir Interview Process (Benchmark)}

\begin{center}
\rowcolors{2}{white}{mslightgray}
\begin{tabular}{C{1cm}L{3.5cm}L{8cm}}
\toprule
\rowcolor{msblue}
\textcolor{white}{\textbf{Stage}} &
\textcolor{white}{\textbf{Format}} &
\textcolor{white}{\textbf{Focus}} \\
\midrule
1 & Online Assessment & Algorithmic problem-solving screen \\
2 & Phone Screen & Technical communication and fit \\
3a & Technical Round 1 & Coding / LeetCode-style problems \\
3b & Technical Round 2 & System Design at scale \\
3c & Technical Round 3 & Problem Decomposition (ambiguous prompts) \\
4 & Hiring Manager & Behavioral, 30-60-90 day plan, culture fit \\
\bottomrule
\end{tabular}
\end{center}

\vspace{4pt}
\textbf{Process metrics} (Glassdoor, 134 interviews): Average duration
28 days. Difficulty: 3.4/5. Positive experience: 61\%.

\textbf{Interview panel}: Technical interviewer, client-facing
interviewer, behavioral interviewer, hiring manager. The multi-faceted
panel reflects the multi-dimensional nature of the role.

\subsection{Common Hiring Mistakes}

\begin{itemize}
    \item \textbf{Over-emphasis on coding}: LeetCode performance is
        necessary but not sufficient. The best coders often make the
        worst FDEs if they lack client-facing skills.
    \item \textbf{Not testing for ambiguity tolerance}: FDEs work in
        environments where requirements are unclear, data is messy, and
        stakeholders disagree. Interviews must simulate this.
    \item \textbf{Ignoring soft skills}: Customer fluency, empathy, and
        business acumen are the true differentiators. A structured
        client-facing simulation should be part of every FDE interview.
    \item \textbf{Skipping the 30-60-90 day plan}: Asking candidates to
        articulate their approach to the first 90 days of a deployment
        reveals strategic thinking and self-direction.
\end{itemize}

% ============================================
% SECTION 9: FUTURE EVOLUTION
% ============================================
\section{Future Evolution: From FDE to Agent Deployment Engineer}

\subsection{The Agentic AI Transition}

As AI systems evolve from tools that assist humans to autonomous agents
that execute tasks independently, the FDE role is undergoing a parallel
transformation. The emerging role of \textbf{Agent Deployment Engineer}
reflects this shift:

\begin{itemize}
    \item \textbf{Traditional FDE}: Builds custom data integrations,
        writes application code, creates workflows on customer
        infrastructure.
    \item \textbf{Agent Deployment Engineer}: Deploys, configures,
        monitors, and governs autonomous AI agents operating within
        customer environments.
\end{itemize}

Celonis's AgentC suite represents the leading edge of this evolution,
with roles explicitly designed around agent lifecycle management rather
than traditional software deployment.

\subsection{Industry Endorsement}

Marty Cagan (Silicon Valley Product Group), one of the most influential
voices in product management, describes himself as ``exceptionally happy''
about the spread of FDE practices. His advice: \textit{``If you do
nothing else, consider having your best engineers visit customer
offices.''}

Gergely Orosz (The Pragmatic Engineer) has published extensive analysis
of the FDE trend, and a16z has endorsed it as a structural shift in
how enterprise software companies organize engineering talent.

\subsection{Career Impact}

The data on FDE alumni careers is striking. Across Palantir, OpenAI,
and other early adopters, FDE alumni have ``disproportionately exceptional
careers'' (multiple sources). The role accelerates career development by
compressing exposure to diverse industries, customers, and technical
challenges into a timeframe that traditional engineering roles cannot
match.

% ============================================
% KPI DASHBOARD
% ============================================
\section{KPI Dashboard}

\begin{center}
\rowcolors{2}{white}{mslightgray}
\small
\begin{tabular}{L{5cm}L{3.5cm}C{1.5cm}L{3cm}}
\toprule
\rowcolor{msblue}
\textcolor{white}{\textbf{KPI}} &
\textcolor{white}{\textbf{Current Value}} &
\textcolor{white}{\textbf{Trend}} &
\textcolor{white}{\textbf{Source}} \\
\midrule
Job posting growth (YoY) &
    +1,165\% & \kpiup & Lightcast \\
Active postings &
    922 (Nov 2025) & \kpiup & Lightcast \\
Median base salary &
    \$173,816 & \kpiup & Bloomberry (1K jobs) \\
Talent supply ratio &
    3:1 (demand:supply) & \kpiup & Addison Group \\
Salary growth rate &
    4.2\% avg, 10\% peak & \kpiup & Pave / industry \\
Top-end TC (Palantir) &
    \$600K+ & \kpiup & Levels.fyi / reports \\
Equity mention rate &
    70\% of postings & \kpiflat & Bloomberry \\
AI Agent skill demand &
    35\% of postings & \kpiup & Bloomberry \\
LLM skill demand &
    31\% of postings & \kpiup & Bloomberry \\
Python requirement &
    66\% of postings & \kpiflat & Bloomberry \\
On-site expectation &
    25--50\% (3--4 days) & \kpiflat & Multiple \\
Burnout indicator &
    2.7/5 WLB (Palantir) & \kpidown & Glassdoor \\
Premium vs SWE &
    25--40\% & \kpiflat & Multiple \\
Startup-size employers &
    58\% of postings & \kpiflat & Bloomberry \\
\bottomrule
\end{tabular}
\end{center}
